%%
%% Author: shoupu_wan
%% 7/19/18
%%

\chapter{Divide and Conquer}\label{ch:divideAndConquer}

\section{Binary Search}\label{sec:binary-search}

%TODO add these topics
Binary search belongs to the divide and conquer algorithm category.

\label{Evernote:Print two binary search trees in ascending order}

It may sound humiliating to a professional programmer to consider or reconsider binary search at
this moment in this chapter.
But reality is that less than $10\%$ of those tested by Jon Bentley can implement binary search
correctly (Jon Bentley).
There are quite a few questions to answer when implementing

When to terminate?
How to handle not-found case?
Where does the program terminate when found?
Which half to turn to and under what condition?
How to use binary search in combination with other tasks?
How to adapt binary search for more complex tasks?

What about equal values?
The treatment of equal values in binary search is subtle and can be very tricky.
All the options are coupled together.

If the symmetry is broken, then it need to broken in a consistent way.
If you leave equal elements on the left partition, then you should compensate for this by
shortening the partition to the left, and vice versa.

Problems with binary search
\begin{enumerate}
    \item Binary search may be equally implemented recursively or iteratively.
    \item refrain from extra 'shortcut' statements other than base case(s).
    For example,
    if m >= e:
    return -1
    \item Pay attention to the flooring/ceiling when obtaining index for the ``middle'' element.
    \item About the exclusiveness/inclusiveness of boundaries: when the search involves searching
    empty ranges or \emph{a} result is not guaranteed, use exclusive boundary `$\left[..\right)$';
    otherwise, use inclusive boundary `$\left[..\right]$'.
\end{enumerate}

\begin{verbatim}
    int* part(int* s, int*e) {
    int* const pivot = s;
    swap(s, s + rand() % (e - s));
    while(true) {
    // the left partition is strictly less than the pivot
    while(s < e && *++s < *pivot); // find the next element >= pivot
    // the right partition is greater than or equal to the pivot
    while(s < e && *--e >= *pivot); // find the prev element < pivot
    if (e <= s) {
    swap(pivot, s - 1); // insert the pivot element
    return s - 1; // shortening the left partition by shrink back the pivot element
    }
    swap(s, e); //swap the out-of-place elements
    }
    }

    int* part1(int* s, int*e) {
    int * p = s + rand() % (e - s);
    swap(s, p);
    p = s;
    for(int* i = s; i < e; i++ ) {
    if (*i < *s) {
    swap(++p, i); //swap the out-of-place elements
    }
    }
    swap(s, p);
    return p;
    }

    void qSort(int* s, int* e) {
    if (e - s < 2) return;
    int* p = part(s, e);
    qSort(s, p);
    qSort(p + 1, e);
    }
\end{verbatim}
Quick sort and binary search are closely related

Binary search may be implemented highly structured

% to be included java example
Example: \code{org.shoupu.binarySearch.SearchRange}

\begin{itemize}
    \item Binary search structure \label{Evernote:Closest Binary Search Tree Value II}
    \item Reusable binary search functions ~\fullref{Evernote:Search for start and end of a target in
    a sorted array}
    \item Post-processing in binary search \label{Evernote:Find median of two sorted arrays}
\end{itemize}

\section{Generic Binary Search Strategy}\label{generic-binary-search}

A lot of the buggy implementations of binary search came from the fitting and nudging work
required to make binary search work for a specific application.
These include but are not limited to infinite loop, out of boundary, skewed choices, and
miss-by-one problem.

I believe making binary search work on the application level is over-engineering.
Many times a lot of work is done to make every bit work together perfectly and this is not
necessary.
In this section, I am going to propose a generic binary search strategy that requires less
implementation effort and yet brings about most of the savings of a precise binary search algorithm
would.
It can be used for many scenarios where binary search is used.
\begin{clisting}
    \lstinputlisting[label={lst:generic-binary-search},
    caption={Structure of a generic binary search.
    Each iteration, it shrinks the search range by half.
    When there are only two choices, we switch to iterative solution instead.
    It is essentially a $80:20$ strategy---with $20\%$ coding effort we achieve $80\%$ of
    the benefit of binary search.}]
    {generic_binary_search.py}
\end{clisting}


\section{Case Study-Hadoop}
As a real-world case, let's dive into the depth of \code{Hadoop} source code and take a look at
the two utility methods in package \ilcode{org.apache.hadoop.io.file.tfile.Utils}
\ilcode{\#lowerBound(..)} and \ilcode{\#upperBound(..)} \footnote{\url{https://github
.com/apache/hadoop/blob/a0da1ec01051108b77f86799dd5e97563b2a3962/hadoop-common-project/hadoop
-common/src/main/java/org/apache/hadoop/io/file/tfile/Utils.java}}.

This is an example where the two binary search implementations could be reduced to one with
little bit structural change.
